% GENERATED FILE. Edit canonical lessons, then run npm run generate:books.
% canonical-source-hash: fnv1a64:b40a634d5a8dd131
% canonical-lessons: KA-C15-niiru-akki, KA-S124-letter-pa, KA-S146-digit-seven

\chapter{Water and Rice}
\label{ch:ka-water-and-rice}
\hlchaptermodality{15}

\begin{chapteropening}
\textbf{By the end of this chapter:} \emph{I can ask for water in Kannada and tell raw rice apart from cooked rice.}

Say \kn{ನೀರು} — matching Tamil's neer, not Malayalam's — then separate \kn{ಅಕ್ಕಿ}, raw rice, from \kn{ಅನ್ನ}, the cooked rice on the plate.
\end{chapteropening}

\section[nīru akki anna]{\kn{ನೀರು}, \kn{ಅಕ್ಕಿ}, \kn{ಅನ್ನ} — water matching Tamil, and rice, raw and cooked}

\label{lesson:KA-C15-niiru-akki}



\begin{quote}
Kannada's water-word sides with Tamil, not with Malayalam's surprising break from the pack.
\end{quote}

\begin{cousinweb}
\textbf{\kn{ನೀರು}} (\emph{nīru}) = \textquotedblleft{}\textbf{water}\textquotedblright{} — matching Tamil's \textbf{\ta{நீர்}} (\emph{nīr}, the root inside \emph{taṇṇīr}) closely. Kannada didn't make Malayalam's switch to the \emph{veḷ-} \textquotedblleft{}bright/white\textquotedblright{} root.
\end{cousinweb}

\begin{cousinweb}
\begin{itemize}
\raggedright
  \item \textbf{\kn{ಅಕ್ಕಿ}} (\emph{akki}) = \textquotedblleft{}\textbf{rice}\textquotedblright{} (raw grain).
  \item \textbf{\kn{ಅನ್ನ}} (\emph{anna}) = \textquotedblleft{}\textbf{cooked rice}\textquotedblright{} (the meal) — likely Sanskrit- influenced, the same general idea as Tamil's \emph{sātam}.
\end{itemize}
\end{cousinweb}

\subsection*{Your turn}
Say these aloud:

\begin{itemize}
\raggedright
  \item \textquotedblleft{}nīru\textquotedblright{} — water, matching Tamil
  \item \textquotedblleft{}akki\textquotedblright{} — raw rice
  \item \textquotedblleft{}anna\textquotedblright{} — cooked rice
\end{itemize}

\begin{tcolorbox}[breakable,colback=teal!4,colframe=teal!35!black,title={Before you move on}]
Does Kannada's water-word match Tamil's or Malayalam's? (\textbf{Tamil's} — \emph{nīru}/\emph{nīr}, not Malayalam's \emph{veḷḷam}.) What's the raw/cooked rice pair? (\textbf{\kn{ಅಕ್ಕಿ}} \emph{akki} raw, \textbf{\kn{ಅನ್ನ}} \emph{anna} cooked.)
\end{tcolorbox}
\section[pa]{\kn{ಪ} — one character, met inside words you already say}

\label{lesson:KA-S124-letter-pa}



\begin{quote}
Before the new one: \kn{◌ಂ} — what does it do?

One character this time. Just one — and you have been saying it for pages without knowing which mark on the page it was.
\end{quote}

\subsection*{Script you'll notice: \kn{ಪ}}
\textbf{\kn{ಪ}} — \emph{pa}.

It is a \textbf{consonant}, and in this script a consonant is never bare: it comes with an \emph{a} already in it. So it is not \emph{p}, it is \textbf{pa}.

You already say these, and every one of them has \kn{ಪ} somewhere inside it:

\begin{itemize}
\raggedright
  \item \textbf{\kn{ಪರವಾಗಿಲ್ಲ}} \emph{paravāgilla} — it's okay / no problem / you're welcome
  \item \textbf{\kn{ಕಪ್ಪು} \kn{ಬಿಳಿ} \kn{ಕೆಂಪು} \kn{ನೀಲಿ}} — black, white, red, blue
  \item \textbf{\kn{ಅಪ್ಪ} \kn{ಅಮ್ಮ} \kn{ಅಣ್ಣ} \kn{ತಮ್ಮ} \kn{ಅಕ್ಕ} \kn{ತಂಗಿ}} — father, mother, and four age-graded sibling words
\end{itemize}

\subsection*{Writing: \kn{ಪ} — copy what you see}
Put your pen on \kn{ಪ} and follow its line. Copy the shape you can see — slowly, and larger than it is printed.

\begin{quote}
This book does not yet tell you \textbf{where to start the character or which way to travel}. That is a real thing, taught with real variation from school to school, and it is not written down here until it can be written down with a source. Copying what is in front of you needs no such source, and it is how the shape gets into your hand in the meantime.
\end{quote}

\subsection*{Your turn}
\begin{itemize}
\raggedright
  \item \emph{Look:} at these words, and find \kn{ಪ} in the ones that have it
\end{itemize}

\begin{quote}
\kn{ಪರವಾಗಿಲ್ಲ}  ·  \kn{ನಮಸ್ಕಾರ}
\end{quote}

\begin{itemize}
\raggedright
  \item \emph{Trace it:} \kn{ಪ} three times, saying \emph{pa} as you finish each one
  \item \emph{Look:} back at any page of this chapter and find \kn{ಪ} once more
\end{itemize}

\begin{tcolorbox}[breakable,colback=teal!4,colframe=teal!35!black,title={Before you move on}]
Which character is this — \kn{ಪ}? What sound does it carry? (\emph{\textbf{pa}}.) Name one word you already say that contains it.
\end{tcolorbox}
\section[7]{\kn{೭} — the digit for ēḷu}

\label{lesson:KA-S146-digit-seven}



\begin{quote}
Before the new shape: \textbf{\kn{೬}} — which number is it?

One digit this time, and you already know its name.
\end{quote}

\subsection*{Script you'll notice: \kn{೭}}
\textbf{\kn{೭}} — \textbf{7}, the digit for \textbf{\kn{ಏಳು}} (\emph{ēḷu}), \textquotedblleft{}seven\textquotedblright{}.

Seven of ten. From here the remaining shapes are the ones you will meet least often, which is exactly why they come last.

\subsection*{Writing: \kn{೭} — copy what you see}
Put your pen on \kn{೭} and follow its line. Copy the shape in front of you — slowly, and larger than it is printed.

\begin{quote}
This book does not yet tell you \textbf{where to start the character or which way to travel}. That is taught with real variation from school to school, and it is not written down here until it can be written down with a source. Copying what you can see needs no such source.
\end{quote}

\subsection*{Your turn}
\begin{itemize}
\raggedright
  \item \emph{Look:} at the digits you have met, and name each one aloud
\end{itemize}

\begin{quote}
\kn{೧}  ·  \kn{೨}  ·  \kn{೩}  ·  \kn{೪}  ·  \kn{೫}  ·  \kn{೬}  ·  \kn{೭}
\end{quote}

\begin{itemize}
\raggedright
  \item \emph{Say it:} \textquotedblleft{}ēḷu\textquotedblright{} as you point at \kn{೭}
\end{itemize}

\begin{tcolorbox}[breakable,colback=teal!4,colframe=teal!35!black,title={Before you move on}]
Which number is \textbf{\kn{೭}}? (\textbf{7} — \emph{ēḷu}, \kn{ಏಳು}.) Name every digit you have met so far, in order.
\end{tcolorbox}
